\uuid{AANr}
\exo7id{3496}
\titre{exo7 3496}
\auteur{quercia}
\organisation{exo7}
\datecreate{2010-03-10}
\isIndication{false}
\isCorrection{true}
\chapitre{Réduction d'endomorphisme, polynôme annulateur}
\sousChapitre{Diagonalisation}
\cpge{mpsi-mp2i.calcul-matriciel-systemes.e, pcsi.determinants.c, ptsi.determinants.a, mpsi-mp2i.determinants.d, pcsi.calcul-matriciel-systemes-lineaires.d, ptsi.calcul-matriciel-systemes-lineaires.d, bcpst1.matrices, mp-mpi.reduction-endomorphismes.b, pc.reduction.a, psi.reduction-endomorphismes.a, pt.reduction-endomorphismes.a, mp-mpi.reduction-endomorphismes.c, pc.reduction.b, psi.reduction-endomorphismes.b, pt.reduction-endomorphismes.b, mp-mpi.reduction-endomorphismes.d, pc.reduction.c, psi.reduction-endomorphismes.c, pt.reduction-endomorphismes.c, bcpst2.valeurs-propres-vecteurs-propres.a, bcpst2.valeurs-propres-vecteurs-propres.b}
\module{Algèbre}
\niveau{L2}
\difficulte{}

\contenu{
\texte{
Diagonaliser les matrices suivantes :
}
\begin{enumerate}
    \item \question{$A=\begin{pmatrix} 1 & 5 \cr 2 & 4 \cr\end{pmatrix}$}
\reponse{$     P=\begin{pmatrix} 1 &-5 \cr 1 & 2 \cr\end{pmatrix}$,\,
$    D= \begin{pmatrix} 6 & 0 \cr 0&-1 \cr\end{pmatrix}$}
    \item \question{$A=\begin{pmatrix} 2 & 5 \cr  4 & 3 \cr\end{pmatrix}$}
\reponse{$    P= \begin{pmatrix} 5 & 1 \cr -4 & 1 \cr\end{pmatrix}$,\,
$     D=\begin{pmatrix}-2 & 0 \cr 0& 7 \cr\end{pmatrix}$}
    \item \question{$A=\begin{pmatrix} 5 & 3 \cr -8 &-6 \cr\end{pmatrix}$}
\reponse{$ P=\begin{pmatrix} 3 & 1 \cr -8 &-1 \cr\end{pmatrix}$,\,
      $ \begin{pmatrix}-3 & 0 \cr 0& 2 \cr\end{pmatrix}$}
    \item \question{$A=\begin{pmatrix} 4 & 4 \cr 1 & 4 \cr\end{pmatrix}$}
\reponse{$P=\begin{pmatrix} 2 &-2 \cr 1 & 1 \cr\end{pmatrix}$,\,
    $ D=\begin{pmatrix} 6 & 0 \cr 0 & 2 \cr\end{pmatrix}$}
\end{enumerate}
}
